\bibitem{refKuchuganov}
Кучуганов, А. В. Методология семантического анализа и поиска графической информации : автореферат диссертации на соискание ученой степени доктора технических наук : специальность 05.13.01 «Системный анализ, управление и обработка информации (в промышленности)» / А. В. Кучуганов. – Ижевск, 2018. – Текст : непосредственный.
\bibitem{refDavletov}
Давлетов, А. Р. Современные методы машинного обучения и технология OCR для автоматизации обработки документов / А. Р. Давлетов. – Текст : электронный // Вестник науки. – 2023. – № 10 (67). – URL: https://cyberleninka.ru/article/n/sovremennye-metody-mashinnogo-obucheniya-i-tehnologiya-ocr-dlya-avtomatizatsii-obrabotki-dokumentov (дата обращения: 22.05.2026).
\bibitem{refRakov}
Раков, Н. С. Машинное зрение / Н. С. Раков, С. В. Пальмов. – Текст : электронный // Форум молодых ученых. – 2018. – № 4 (20). – URL: https://cyberleninka.ru/article/n/mashinnoe-zrenie (дата обращения: 22.05.2026).
\bibitem{refRublevsky}
Рублевский, Р. Ландшафт основных терминов в области генеративного AI, их взаимосвязь и употребление / Р. Рублевский. – Текст : электронный // Хабр. – 2025. – 25 сент. – URL: https://habr.com/ru/articles/950600/ (дата обращения: 22.05.2026).
\bibitem{refAkhil}
Akhil, S. Overview of Tesseract OCR engine / S. Akhil. – Текст : электронный // ResearchGate. – 2016. – Dec. – URL: https://www.researchgate.net/publication/315834331 (дата обращения: 22.05.2026).